\chapter*{术语表}
\addcontentsline{toc}{chapter}{术语表}

\begin{description}
  \item[RAII] Resource Acquisition Is Initialization。把资源生命周期绑定到对象生命周期。
  \item[ABI] Application Binary Interface，应用二进制接口，规定调用约定、对象布局、符号和链接边界。
  \item[API] Application Programming Interface，源码层接口，规定调用者能传什么、得到什么、错误怎样表达。
  \item[ID] Identifier，身份或编号。业务代码中常用它唯一定位用户、玩家、记录或资源。
  \item[C++14/C++20/C++23] C++ 标准版本名。C++14、C++20、C++23 表示语言和标准库能力逐步增加；工程中必须按目标编译器和平台实际支持选择。
  \item[UB] Undefined Behavior，未定义行为。触发后标准不再规定程序结果。
  \item[I/O] Input/Output，输入输出。C++ 中常通过 stream、文件描述符或库接口表达。
  \item[EOF] End Of File，输入结束。它不是一个普通字符，而是流状态或读取结果的一部分。
  \item[ASCII] 早期单字节字符集合，本书简单配置和英文命令示例多按它处理。
  \item[Unicode] 统一字符集合。中文、emoji、组合字符等场景需要专门文本处理边界。
  \item[UTF-8] Unicode 的常见变长字节编码；\code{std::string::size()} 返回字节数，不是用户感知字符数。
  \item[IDE] Integrated Development Environment，集成开发环境，通常包含编辑器、构建入口、调试器和代码导航。
  \item[GUI] Graphical User Interface，图形用户界面。与 CLI 相对，通过窗口、按钮和菜单交互。
  \item[CI] Continuous Integration，持续集成。自动构建、测试和检查每次提交。
  \item[CLI] Command-Line Interface，命令行接口。
  \item[CMake] C++ 常用构建系统，用 target 描述可执行文件、库、依赖和编译选项。
  \item[GNU] GNU 工具链生态名，\code{g++}、\code{ld}、\code{gdb} 等工具常来自这个生态。
  \item[HTTP] HyperText Transfer Protocol，超文本传输协议，常用于服务 API、测试桩和外部适配器。
  \item[\code{std::}] C++ 标准库命名空间前缀。看到 \code{std::vector} 先把它理解成“拥有元素的动态数组”，看到 \code{std::span} 先理解成“不拥有元素的连续视图”。
  \item[\code{size\_t}] C++ 容器大小和下标常用类型，通常是无符号整数。和负数或倒序循环混用时要特别谨慎。
  \item[\code{int64\_t/uint64\_t}] 固定 64 位整数类型。有符号版本能表达负数，无符号版本更适合位模式、哈希和容量上界。
  \item[target] CMake 中的构建目标，例如库、可执行文件或测试程序。
  \item[PUBLIC/PRIVATE/INTERFACE] CMake 依赖传播关键字。PUBLIC 影响当前 target 和使用者，PRIVATE 只影响当前 target，INTERFACE 只表达使用要求。
  \item[\code{DCMAKE\_BUILD\_TYPE}] CMake 单配置生成器常用变量，用来选择 Debug、Release、RelWithDebInfo 等构建类型。
  \item[\code{DCMAKE\_CXX\_FLAGS}] CMake 中追加 C++ 编译参数的变量。工程中更推荐优先使用 target 级别的编译选项，减少全局副作用。
  \item[命名常量] 全大写的业务或工程边界名，例如 \code{PORT_MIN}、\code{RETRY_MAX_ATTEMPTS}。它们的价值是把数字背后的约束说清楚。
  \item[MB] Megabyte，兆字节量级。性能和内存章节用它描述数据规模或分配量。
  \item[ASan] AddressSanitizer，地址消毒器，用于发现越界、use-after-free 等内存错误。
  \item[UBSan] UndefinedBehaviorSanitizer，用于发现有符号溢出、非法移位、错误对齐等未定义行为。
  \item[TSan] ThreadSanitizer，用于发现数据竞争和部分同步错误。
  \item[TLS] Thread-Local Storage，线程局部存储，每个线程拥有一份变量实例。
  \item[JSON] JavaScript Object Notation，一种常用文本数据格式。
  \item[CPU] Central Processing Unit，中央处理器。
  \item[O0/O2] 编译优化级别。O0 便于调试，O2 代表常用优化组合，二者生成代码和调试体验会明显不同。
  \item[DNDEBUG] 预处理宏，常用于关闭 \code{assert}。发布构建中打开它后，依赖断言的副作用会消失。
  \item[LRU] Least Recently Used，最近最少使用淘汰策略。
  \item[FIFO] First In First Out，先进先出顺序，常用于队列。
  \item[不变量] 对象在对外可见状态下必须始终成立的条件。
  \item[值语义] 对象可以像普通值一样复制、赋值和组合，副本之间互不影响。
  \item[所有权] 谁负责释放资源、谁能延长资源生命周期的规则。
  \item[悬空引用] 引用或指针指向的对象已经销毁。
  \item[迭代器失效] 容器操作后，原迭代器不再安全可用。
  \item[concept] C++20 中用于约束模板参数能力的语言机制。
  \item[数据竞争] 多个线程并发访问同一对象，至少一个写入，且没有同步。
  \item[视图] 观察已有数据的轻量对象，通常不拥有底层数据。
  \item[严格弱序] 排序比较器需要满足的关系，相等元素之间不能互相小于对方。
  \item[死锁] 多个执行单元互相等待对方释放资源，导致都无法继续。
  \item[RAII 锁] 用对象构造和析构管理加锁、解锁的同步写法。
  \item[lambda] 可以在局部定义的匿名可调用对象，常用于谓词、回调和小策略。
  \item[类型擦除] 隐藏具体可调用对象或实现类型，只暴露统一接口的技术。
  \item[异常安全] 异常发生后对象和程序状态仍满足约定保证的能力。
  \item[强保证] 函数失败后状态像没有调用过一样。
  \item[虚函数] 通过基类接口在运行期分派到派生类实现的机制。
  \item[variant] 表示若干已知类型之一的标准库类型，适合封闭类型集合。
  \item[背压] 下游处理不过来时，让上游减速、失败或丢弃任务的过载控制机制。
  \item[对齐] 对象在内存中按平台要求放置到特定地址边界的规则。
  \item[填充字节] 编译器为满足对齐或对象布局要求插入的无业务含义字节。
  \item[特征测试] 在重构前固定现有外部行为的测试，用于防止行为漂移。
  \item[重载] 多个函数共享同一名字，但参数列表不同，由编译器根据调用选择。
  \item[命名空间] 控制名字归属和范围的语言机制，用于组织模块并减少冲突。
  \item[词汇类型] 标准库中表达常见语义的类型，例如 \code{optional}、\code{variant}、\code{span}。
  \item[顺序一致] 默认原子内存序，提供最容易推理的全局一致原子操作顺序。
  \item[逻辑竞争] 没有数据竞争但业务检查和更新不是一个原子事务，导致结果违反业务规则。
  \item[constexpr] 表示表达式或函数可用于编译期求值的语言机制。
  \item[type traits] 在编译期查询类型属性的一组模板工具。
  \item[SFINAE] 模板替换失败不构成立即编译错误的规则，常用于旧式模板约束。C++20 concepts 能把许多 SFINAE 条件写得更直接。
  \item[非类型模板参数] 模板参数不是类型而是值，例如 \code{RingBuffer<T, CAPACITY>} 中的容量。
  \item[协程] 可以暂停和恢复的函数形态，用于表达异步控制流。
  \item[协程帧] 协程暂停时保存局部状态和恢复位置的存储区域。
\end{description}
